% Generated draft: requires mathematical and source review.
\documentclass[11pt]{article}
\usepackage[margin=25mm]{geometry}
\usepackage{fontspec}
\setmainfont{Latin Modern Roman}
\usepackage{amsmath,amssymb,mathtools,unicode-math}
\setmathfont{Latin Modern Math}
\usepackage{enumitem,needspace,xurl,hyperref}
\setlength{\parindent}{0pt}
\setlength{\parskip}{5pt}
\setlength{\emergencystretch}{3em}
\newcommand{\N}{\mathbb{N}}
\newcommand{\Z}{\mathbb{Z}}
\newcommand{\Q}{\mathbb{Q}}
\newcommand{\R}{\mathbb{R}}
\newcommand{\C}{\mathbb{C}}
\newcommand{\F}{\mathbb{F}}
\newcommand{\A}{\mathbb{A}}
\newcommand{\PP}{\mathbb{P}}
\newcommand{\E}{\mathbb{E}}
\newcommand{\eps}{\varepsilon}
\newcommand{\dd}{\,\mathrm{d}}
\begin{document}
\section*{UnsolvedMath ID: 9700027}
\textbf{Source problem code:} catalogue entry
\subsection{Definitions and mathematical statement}
In the usual planar spatial-route setting, a random route $\mathcal R(z,z')$ joins every pair of points of $\mathbb R^2$, with joint consistency/compatibility. For distinct $z_1,\ldots,z_k$, its finite-point span is the embedded network
\[
\operatorname{span}(z_1,\ldots,z_k)=\bigcup_{1\leq i<j\leq k}\mathcal R(z_i,z_j).
\]
Similarity covariance means that for translations, rotations, and dilations $S(z)=aRz+b$, $a>0$, the joint law of transformed routes equals the law of routes between the transformed endpoints. A scale-invariant random spatial network (SIRSN) also has route-compatibility and quantitative finiteness/integrability axioms; the supplied item does not reproduce those axioms.

Construct a SIRSN whose law is mathematically natural. The source offers an explicit formula for the law of $\operatorname{span}(z_1,\ldots,z_k)$ as an example of such naturality, rather than the sole stated criterion. Such a formula should work for arbitrary finite point configurations and specify the embedded-network distribution, not merely a mean length or one marginal statistic.

The supplied summary identifies an Aldous--Kendall improper-Poisson-line construction as a natural candidate and says explicit span distributions remain unavailable. That report does not turn the optional example into a mandatory source requirement. The meaning of mathematically natural and the representation expected of an explicit law are not fixed, so the task remains a construction/tractability program with the full SIRSN conditions to be checked.
\subsection{Short English statement}
Construct a mathematically natural scale-invariant random spatial network, ideally with explicit finite-point span laws. The exact SIRSN axioms and criterion of naturality are not reproduced, and the explicit-law example is optional in the original wording.
\subsection{Sources}
\begin{itemize}
\item David Aldous, \emph{Open problems}, page identified as \texttt{sirsn.html}, linked paper Open Problem 27, maintained index \url{https://www.stat.berkeley.edu/~aldous/Research/OP/index.html}; accessed 2026-07-29. The supplied extraction includes HTML and linked primary PDF text but no separate PDF identifier.
\item Aldous--Kendall improper Poisson line construction, as the candidate-model attribution in the supplied research summary; no complete publication citation is given.
\item Local registry \texttt{data/unsolvedmath/problems.json}, permanent UnsolvedMath ID $9700027$, source code catalogue entry, \emph{A mathematically natural SIRSN}. Dataset attribution: Ulam AI and the UnsolvedMath Contributors.
\end{itemize}
\end{document}
